% ---------------------------------------------------------------------------
% Methods: rule distillation (NtRED) and its evaluation.
%
% Draft for review, 2026-08-08. Replaces the v1/v2 description in
% docs/DISTILLATION_V2.md. Written to be implementable: every quantity is
% defined, every hyperparameter is named with the value we measured, and the
% objective is stated explicitly.
%
% The substantive change from v1/v2 is in Section "Alignment": the previous
% matcher decided both what counts as a miss (the distillation TARGET) and what
% counts as a false positive (the distillation GATE) using a text-only,
% one-to-one, per-type rule. Measured on 61 held-out papers, that rule mislabels
% 39% of false positives and 21% of misses. Since the same rule defines the
% objective, the loop was optimising against a wrong target, not merely
% reporting one.
% ---------------------------------------------------------------------------

\section{Methods}
\label{sec:methods}

\subsection{Setting and notation}
\label{sec:notation}

A paper $P$ is processed by an extractor into a set of typed statements. Write
$\mathcal{T}$ for the node-type vocabulary (Section~\ref{sec:schema};
$|\mathcal{T}| = 18$), and represent an extracted item as a pair
$e = (\tau_e, x_e)$ with $\tau_e \in \mathcal{T}$ its type and $x_e$ its
verbatim text span.

Two extractors are compared throughout:

\begin{itemize}
  \item $\mathcal{L}$, the \emph{teacher}: a large language model prompted with
    the schema, producing $\mathcal{L}(P)$. Its output is the silver standard.
  \item $\mathcal{R}_\Theta$, the \emph{student}: a deterministic rule engine
    parameterised by an ordered set of regular-expression rules $\Theta$,
    producing $\mathcal{R}_\Theta(P)$.
\end{itemize}

The distillation problem is to search for $\Theta$ such that
$\mathcal{R}_\Theta$ reproduces $\mathcal{L}$ on unseen papers, subject to
$\Theta$ remaining small and human-readable. The interpretability constraint is
not decoration: a rule set that reproduced $\mathcal{L}$ perfectly but was
unreadable would restate the teacher rather than explain it.

We write $D_{\mathrm{train}}$ and $D_{\mathrm{test}}$ for disjoint paper sets,
and $G(P) = \mathcal{L}(P)$ for the reference (``gold'') items of paper $P$.

\subsection{Alignment between two extractions}
\label{sec:alignment}

Every quantity below --- the training signal, the rule statistics, and the
reported scores --- rests on deciding when an extracted item corresponds to a
reference item. This is stated first because it is the component that changed.

\paragraph{Item similarity.}
For items $e$ and $g$ define
\begin{equation}
  s(e, g) \;=\; \sigma\!\left(x_e, x_g\right) \cdot \kappa(\tau_e, \tau_g),
  \label{eq:sim}
\end{equation}
where $\sigma$ is a text similarity on normalised token sets and $\kappa$ is a
type-compatibility factor.

For $\sigma$ we use \emph{containment},
$\sigma(x,y) = |T(x) \cap T(y)| \,/\, \min(|T(x)|, |T(y)|)$, with $T(\cdot)$ the
token set after LaTeX normalisation (macro expansion against the paper's own
glossary, environment stripping, case folding; Section~\ref{sec:normalisation}).
Containment rather than Jaccard because a correct extraction that includes a
trailing clause the reference omits should not be penalised for length.

$\kappa$ is a graded matrix over $\mathcal{T} \times \mathcal{T}$ with
$\kappa(\tau,\tau) = 1$, $\kappa$ large for near-synonymous types
(\texttt{theorem}/\texttt{proposition}), small for structurally distinct ones
(\texttt{definition}/\texttt{proof}), and $0$ where a match is meaningless. The
matrix is given in Appendix~\ref{app:typecompat} and is shared with the
cross-paper alignment task, so a single definition of ``these two statements are
the same thing'' serves both.

\paragraph{Group alignment.}
A one-to-one matching cannot express the two commonest disagreements between
extractors: a reference statement emitted as two items (\emph{split}), and two
reference statements emitted as one (\emph{merge}). Both are presentational, not
semantic. We therefore align \emph{groups}.

Given $E = \mathcal{R}_\Theta(P)$ and $G$, construct the alignment $A(E,G)$:

\begin{enumerate}
  \item \textbf{Unambiguous pairs first.} Form all pairs with
    $s(e,g) \ge \theta$, sort by $s$ descending, and greedily accept a pair when
    neither element is already used. This reduces exactly to the classical
    one-to-one matching when no groups are needed, so the alignment is a strict
    generalisation of the previous rule.
  \item \textbf{Splits.} For each unmatched $g \in G$, let
    $C_g = \{e \text{ unmatched} : \kappa(\tau_e, \tau_g) > 0,\;
      |T(x_e) \cap T(x_g)| \ge m_{\mathrm{tok}},\;
      J(x_e, x_g) \ge m_J \}$
    be its candidate members, capped at the $k_{\max}$ largest by overlap.
    If $\sigma\big(\bigcup_{e \in C_g} x_e,\; x_g\big) \ge \theta$, accept the
    group: mark $g$ and all of $C_g$ matched.
  \item \textbf{Merges.} Symmetrically, for each unmatched $e \in E$ against
    unmatched reference items.
\end{enumerate}

The guards $m_{\mathrm{tok}}$ and $m_J$ are load-bearing rather than
defensive. Without them, pooling enough unrelated text drives containment
against a short target to $1$ and the group rule accepts anything; in an early
version this inflated $F_1$ from $0.52$ to $0.86$ before the guards were
restored. They also matter empirically: with pooling unguarded, measured
precision on our test set read $0.93$; with the guards it is $0.75$.

\paragraph{Item-level counting.}
Under group alignment a ``true positive'' is not a pair --- one group match can
credit three extracted items against two reference items --- so precision and
recall are computed over \emph{items covered}:
\begin{equation}
  \mathrm{Prec} = \frac{|\{e \in E : e \text{ matched}\}|}{|E|},
  \qquad
  \mathrm{Rec} = \frac{|\{g \in G : g \text{ matched}\}|}{|G|},
  \qquad
  F_1 = \frac{2\,\mathrm{Prec}\,\mathrm{Rec}}{\mathrm{Prec} + \mathrm{Rec}}.
  \label{eq:prf}
\end{equation}
When the alignment is one-to-one these reduce to the usual $TP/(TP{+}FP)$ and
$TP/(TP{+}FN)$, so the two regimes are directly comparable.

Scores are \emph{pooled} over $(\text{paper})$ rather than macro-averaged: per
paper and per type, counts are small and unstable, and macro-averaging over
sparse cells lets a paper with two items dominate one with two hundred.

Hyperparameters: $\theta = 0.7$, $m_{\mathrm{tok}} = 5$ tokens, $m_J = 0.10$,
$k_{\max} = 12$.

\subsection{Dark knowledge}
\label{sec:dk}

The teacher's advantage over the student on paper $P$ is the set of reference
items the student fails to cover,
\begin{equation}
  \mathrm{DK}_\Theta(P) \;=\; \{\, g \in G(P) \;:\; g \text{ unmatched in }
    A(\mathcal{R}_\Theta(P),\, G(P)) \,\}.
  \label{eq:dk}
\end{equation}
This is the object the loop tries to eliminate, and the reason the alignment
definition matters so much: under a one-to-one matcher, a reference statement
that the student \emph{did} recover as two items is scored as dark knowledge,
and the loop then synthesises a rule for a pattern already covered. Such a rule
cannot improve recall and can only add false positives --- which is the
mechanism by which precision fell while recall rose in the previous
configuration (Section~\ref{sec:results}).

\subsection{Distillation loop}
\label{sec:loop}

Rules live in a persistent registry $\mathcal{G}$, not in the emitted source
file. Each rule $r$ carries an immutable identifier $\mathrm{uid}(r)$, a regular
expression, the target type it emits, the paper it was born from, a status in
$\{\textsc{active}, \textsc{retired}\}$, and per-paper statistics. The emitted
\texttt{rules.py} is a projection of $\{r : \mathrm{status}(r) =
\textsc{active}\}$. Identifiers are never reused, and rules are never deleted;
only the status flips. This makes gating reversible and makes an $\alpha$-sweep
a dictionary scan rather than a sequence of file snapshots.

\begin{algorithm}
\caption{Distillation over a training stream}
\label{alg:distill}
\begin{algorithmic}[1]
\Require papers $D_{\mathrm{train}}$, teacher outputs $G(\cdot)$, budget $B$
\State $\mathcal{G} \gets \textsc{LoadRegistry}()$
\ForAll{$P \in D_{\mathrm{train}}$}
  \State $E \gets \mathcal{R}_{\mathrm{active}(\mathcal{G})}(P)$
    \Comment{deterministic; no model call}
  \State $\mathrm{DK} \gets \textsc{Unmatched}(A(E, G(P)))$
    \Comment{Eq.~\ref{eq:dk}, group-aware}
  \Statex
  \State \textbf{Recycle.} \textbf{for} $r \in \mathrm{retired}(\mathcal{G})$
    \textbf{do if} $r$ fires on any $d \in \mathrm{DK}$ \textbf{then}
    reactivate $r$; remove the items it covers from $\mathrm{DK}$
    \Comment{no model call}
  \Statex
  \If{$\mathrm{DK} \ne \emptyset$ \textbf{and} budget remains}
    \State $S \gets \textsc{Propose}(\mathrm{DK}, \mathrm{tex}(P))$
      \Comment{model call 1: candidate patterns}
    \State $S' \gets \textsc{Integrate}(\mathrm{active}(\mathcal{G}), S)$
      \Comment{model call 2: merge, do not rewrite}
    \ForAll{$(\text{pattern}, \tau) \in S'$}
      \State $\mathcal{G}.\textsc{MintOrUpdate}(\text{pattern}, \tau,
        \text{born}{=}P)$
    \EndFor
  \EndIf
  \Statex
  \State \textbf{Attribute.} \textbf{for} $r \in \mathcal{G}$: recompute
    $(\mathrm{tp}_r(P), \mathrm{fp}_r(P))$ on every paper seen so far
    \Comment{no model call}
  \State \textbf{Gate.} $\mathrm{status}(r) \gets \textsc{Active}$ \textbf{iff}
    $r$ is retained by the selection rule (Section~\ref{sec:objective})
  \State \textbf{Emit.} write \texttt{rules.py} from
    $\mathrm{active}(\mathcal{G})$; persist $\mathcal{G}$
\EndFor
\end{algorithmic}
\end{algorithm}

Three properties are worth stating because they are what make the loop cheap and
auditable. Only line 8--9 calls a model, so cost is proportional to the number of
papers with non-empty $\mathrm{DK}$ rather than to $|D_{\mathrm{train}}|$; in
practice this falls off sharply as coverage grows. \textsc{Integrate} is given
the current rule set and asked to \emph{add}, never to rewrite, so identifiers
survive and a rule's statistics accumulate across papers instead of resetting.
And because every rule retains its per-paper counts even while retired,
re-gating at a different operating point requires no re-extraction.

\paragraph{Attribution.}
A rule $r$ is credited with $\mathrm{tp}_r(P)$, the number of items it emitted
that are matched in $A(E, G(P))$, and $\mathrm{fp}_r(P)$, the number unmatched.
Where an item is produced by more than one rule, credit is assigned to the
earliest-firing rule, so the counts sum to the item totals and precision is not
double-counted.

\subsection{Selection objective}
\label{sec:objective}

Let $\Theta \subseteq \mathcal{G}$ be a candidate active set and let
$F_1(\Theta; D)$ be Eq.~\ref{eq:prf} evaluated with $\mathcal{R}_\Theta$ on
paper set $D$. We select
\begin{equation}
  \Theta^\star \;=\; \arg\max_{\Theta \subseteq \mathcal{G}}
    \; F_1(\Theta;\, D_{\mathrm{val}}) \;-\; \lambda \, |\Theta| .
  \label{eq:objective}
\end{equation}

Two comments on this choice, since it differs from what was used previously.

\emph{Why $F_1$ and not per-rule false-discovery rate.} The earlier loop retired
any rule whose FDR exceeded a threshold $\alpha$. That optimises precision in
isolation and is blind to what retiring a rule costs in recall: a rule with
mediocre precision that is the \emph{only} rule covering a class of statements is
worth keeping. Eq.~\ref{eq:objective} scores the set, not the rule.

\emph{Why the $\lambda|\Theta|$ term.} Without it the objective is monotone in
rule count --- one can always add a rule that memorises a single paper --- and
the interpretability claim quietly evaporates. $\lambda$ is the price in $F_1$ we
are willing to pay per rule, and reporting the frontier over $\lambda$ is more
honest than reporting a single point, because it exposes the accuracy/complexity
trade-off the paper is actually making.

Exact maximisation is combinatorial. We use backward greedy elimination: start
from $\Theta = \mathcal{G}$, repeatedly retire the single rule whose removal most
increases the objective, and stop when no removal helps. Each step is a lookup
over cached per-paper counts, so a full pass costs no extraction. The result is a
local optimum; we report it as such.

$\lambda$ and $\theta$ are chosen by cross-validation on
$D_{\mathrm{train}}$ only, and $D_{\mathrm{test}}$ is scored once.

\subsection{What is reported}
\label{sec:reporting}

For the distillation claim to be checkable, three things are reported together
and never separately:

\begin{enumerate}
  \item \textbf{Coverage.} $\mathrm{Prec}$, $\mathrm{Rec}$, $F_1$
    (Eq.~\ref{eq:prf}) for the base extractor, the distilled extractor, and the
    teacher, on $D_{\mathrm{test}}$. The base extractor is the honest baseline:
    the gain attributable to distillation is $F_1(\text{distilled}) -
    F_1(\text{base})$, not the absolute number.
  \item \textbf{Complexity.} $|\Theta^\star|$ and total pattern length, with the
    frontier over $\lambda$.
  \item \textbf{Cost.} Model calls and tokens, split into \textsc{Propose} and
    \textsc{Integrate}, and the fraction of papers that required either.
\end{enumerate}

Because the alignment defines the objective, any comparison against previously
published numbers must state which alignment produced them; the same extractor
scores differently under the two, by more than the effect being measured
(Section~\ref{sec:results}).

\subsection{Ablations}
\label{sec:ablations}

The claims above are separable and each is tested by removing one component:
group alignment (revert to one-to-one), type grading (revert $\kappa$ to the
identity), recycling, the complexity penalty ($\lambda = 0$), and
\textsc{Integrate}'s add-only constraint. The alignment ablation is the
important one, since it is the only change that moves the training target rather
than the score.
